\chapter{The Nature of Everything}

\section{Introduction}

This text is a chapter consisting of sections, each consisting of paragraphs, text that consists of 
words, each constiting of minimal primitives called characters, digits and other special symbols.

Ultimately, in digital computer, each one is an integer number, such as $12$ - a mathematical object,
consisting of two digits of power of ten:

\[
1 * 10^1 + 2 * 10^0
\]

But base ten is just a convenient choice for humans with ten fingers.
Each number is an object having ten possible internal states.

But we are free to choose any representation we like, such as a binary base, in which the decimal number $12$
can be represented as a binary number $1010$.

\[
1*2^3 + 0 * 2^2 + 1 * 2^1 + 0 * 2^0
\]

This is all obvious and self evident. 

Integer numbers are just one of the many forms of information. 


\section{What is Information?}

There is no universal definition of information. Instead, there are several well-established definitions,
describing different aspects of the concept. Most commonly used definitions include:

\begin{itemize}
\item \textbf{Shannon}	Reduction of uncertainty, quantified through probability
\item \textbf{Hartley}	Selection from a set of possibilities
\item \textbf{Fisher}	Information about an unknown parameter carried by observations
\item \textbf{Quantum information}	Information represented by quantum states and their distinguishability
\end{itemize}

Most notably, information needs attributes such as distinguishability. Therefore, one might think that
binary object bit would be fundamental minimum.  Bitstring would therefore be fundamental substrate of the informational universe. 


\subsection{Naive Interpretation of Bitstring}

The assumption that information requires distinguishability, therefore binary is the simplest possible representation
suffers from the same problem as the current theories of physics; binary already assumes quite a lot:

\begin{itemize} 
\item Two distinguishable states.
\item An ordering or addressing scheme if the bits form a string.
\item Interpretation of semantics to one or zero (or perhaps to heads and tails)
\item An interpretation that assigns meaning to the resulting arrangement.
\end{itemize}


If we reject all external readers (interpreters), not many of these assumptions remain.

A collection of binary objects without an interpretation is not a bitstring (ordered set of bits).
It is an unordered collection whose only remaining attribute is cardinality - the number of bits. 

In other words:

\[
\underbrace{01001010}\_{\text{interpreted bitstring}} \quad\longrightarrow\quad \underbrace{\{0,1,0,0,1,0,1,0\}}\_{\text{unordered binary objects}}
\]

The right-hand side is still a collection of binary objects, but its ordering and meaning have disappeared.


\section{Unary String}

One might ask, can information be defined with unary objects that can only have one internal state.

Just like one can have a set of bits, one could have a set of unaries (let's call them uits).

If the only available property of the unordered collection is its cardinality, the observer can introduce a distinction between two parts.
In static timeless universe without external metaphysics observer is emrgent, a structure in the information itself.  

But something interesting emerges with the observer: the two binary values are no longer necessary.

The observer–observed distinction is itself sufficient to establish a partition.

Consider an uitstring

|||

Just like binary number in computer does not say how it should be interpreted unsigned value of $3$ or signed $-1$, the above unary string
can be interpreted as:

| ||
|| |


Once we have a partition, we can consider the multiplicity of its constituents.

We can therefore replace the binary objects with indistinguishable unary marks.

Conceptually:

\[
\text{binary objects} \longrightarrow \text{unary multiplicity}
\]

Conclusion: binary distinguishability is not required as a primitive.
It is a possible consequence of an observer interpreting a more primitive multiplicity.




\section{Can we go below unary?}

Can we go even further down the fundamental substrate of everything? 

A unary collection already presupposes:

1. Something that can be counted.
2. A notion of multiplicity.
3. A distinction between one and another, even if the individual elements are indistinguishable.

What remains?

We cannot say that there is one thing, because saying one already invokes cardinality.

We cannot say that there are no things either, because that would introduce a distinction between existence and nonexistence.

We cannot describe it as a set, because a set already requires membership and identity.

We can only assue the fundamental substrate is infinite. If it was finite, then finite by how many bits?  Why $n$, why not $n+1$?

Perhaps the only defensible description is \textit{undifferentiated potentiality}. However, potentiality is already a concept with semantic content.

A true theory of everything cannot exlude any semantics either.

Therefore, we conclude that the obsever is everything, and there is nothing outside it. 

Observer introduces distinquishability into two sub-sets: who observes, and what is being observed.
Observer is the reader, an emergent one, giving meaning, semantics, reading order, and distinquishability. 


Before the split, there is no distinction.

After the split, there is:
\begin{itemize}
\item An observer.
\item An observed part.
\item  A relationship between them.
\item A possible difference across their interface.
\end{itemize}

And from that difference, multiplicity, ordering and eventually binary representation can emerge.

This suggests a potentially powerful reversal of the conventional information-theoretic perspective:

Information does not necessarily require pre-existing distinguishable objects. Distinguishability can emerge from the existence of an observer.

\section{Definition of Obsevation}

Observation is relational property of the observer–observed interface.

When both sides have identical internal structure, the interface is null: nothing can be distinguished across their boundary.

An example of this is a system with camera film and the scene to be filmed.

if the state of the film matches the size of the scene, the camera is blid to the scene. The
interface between the two is null.

